% ===== Document Preamble =====
\documentclass[11pt]{article}
\usepackage[margin=0.5in]{geometry}
\usepackage{amsmath,amssymb}
\usepackage{xcolor}
\usepackage{enumitem}
\usepackage{needspace}
\usepackage{comment}

% ===== Colors =====
\definecolor{AnsRed}{RGB}{170,0,0}
\definecolor{WrongBlue}{RGB}{0,0,160}
\definecolor{SkipGray}{RGB}{120,120,120}
\definecolor{ScenarioTeal}{RGB}{0,100,100}

% ===== Choice list styling (tight, consistent) =====
\setlist[enumerate]{%
  leftmargin=1.5em,
  itemsep=0.25em,
  topsep=0.25em,
  label=\textbf{\Alph*.}
}

% ===== Style Wrappers =====
\newcommand{\QuestionStyle}[1]{{\large\bfseries\boldmath #1}}
\newcommand{\AnswerStyle}[1]{{\large\bfseries\boldmath #1}}
\newcommand{\ReasonStyle}[1]{{\normalsize #1}}

% ===== Show / Hide Reasoning Toggle =====
\newif\ifshowreason
\showreasontrue      % Instructor version
% \showreasonfalse   % Student version

% ===== Choice Item Macros (Answer + Reasoning) =====
% Usage:
%   \Correct{Answer}{Reasoning}
%   \Wrong{Answer}{Reasoning}

\newcommand{\Correct}[2]{%
  \item \AnswerStyle{\textcolor{AnsRed}{#1}}%
  \ifshowreason\\[-0.2em]
    \ReasonStyle{\textcolor{AnsRed}{#2}}%
  \fi
}

\newcommand{\Wrong}[2]{%
  \item \AnswerStyle{\textcolor{WrongBlue}{#1}}%
  \ifshowreason\\[-0.2em]
    \ReasonStyle{\textcolor{WrongBlue}{#2}}%
  \fi
}

% ===== Question Block Macro =====
% Usage: \QBlock{Q1}{stem...}{<items>}

\newcommand{\QBlock}[3]{%
  \Needspace{8\baselineskip}%
  \noindent
  \begin{minipage}{\linewidth}
    \QuestionStyle{#1.~#2}

    \vspace{0.25em}
    \begin{enumerate}
      #3
    \end{enumerate}

    \vspace{0.25em}
    \hrule height 0.4pt
  \end{minipage}
  \vspace{0.75em}
}

% ===== Open-Ended Question Block Macro =====
% Usage: \OEBlock{Q1}{stem...}

\newcommand{\OEBlock}[2]{%
  \Needspace{5\baselineskip}%
  \noindent
  \begin{minipage}{\linewidth}
    \QuestionStyle{#1.~#2}

    \vspace{0.25em}
    \hrule height 0.4pt
  \end{minipage}
  \vspace{0.75em}
}

% ===== Scenario Preamble Block Macro =====
% Usage: \ScenarioBlock{scenario text}

\newcommand{\ScenarioBlock}[1]{%
  \Needspace{4\baselineskip}%
  \noindent
  \begin{minipage}{\linewidth}
    {\normalsize\textcolor{ScenarioTeal}{\textbf{Scenario:} \textit{#1}}}
    \vspace{0.25em}
    \hrule height 0.4pt
  \end{minipage}
  \vspace{0.5em}
}

% ===== Skipped Question Note Macro =====
% Usage: \SkipBlock{Q1}{reason}

\newcommand{\SkipBlock}[2]{%
  \noindent\textcolor{SkipGray}{\textit{[#1 omitted: #2]}}%
  \vspace{0.5em}\par
}

% ===== Document Start =====

\begin{document}

\begin{center}
{\huge \textbf{Quantum Mechanics: Wavefunctions and Stationary States}}\\[0.25em]
{\Large Open-Ended Practice Questions}\\[0.25em]
{\large Questions generated below. Multiple-choice answers in RED; incorrect answers in BLUE.}
\end{center}

\vspace{0.75em}

\OEBlock{Q1}{A particle of mass $m$ in one dimension has the wave function
\[
\Psi(x,t)=A\exp\!\left[-\alpha\left(\frac{m x^2}{\hbar}+ i\,2t\right)\right],
\]
where $\alpha>0$ is a constant.\par
\textbf{(a)} Determine the normalization constant $A$.\par
\textbf{(b)} Write the integral expression used to compute the expectation value $\langle x\rangle$ and evaluate $\langle x\rangle$ for this state.}

\OEBlock{Q2}{At time $t=0$, a particle is represented by the (real) wave function
\[
\Psi(x,0)=\begin{cases}
A\,(x/a), & 0\le x\le a\\
A\,(b-x)/(b-a), & a\le x\le b\\
0, & \text{otherwise}
\end{cases}
\]
where $a$ and $b$ are positive constants with $0<a<b$, and the constant $A$ is chosen so that $\Psi$ is normalized.\par
\textbf{(a)} Find $A$ in terms of $a$ and $b$.\par
\textbf{(b)} What is the probability of finding the particle to the left of $x=a$ (i.e., for $x<a$)?}

\OEBlock{Q3}{Consider a quantum state of the form
\[
\Psi_n(x,t)=\phi_n(x)\,e^{-i\omega_n t},
\]
where $\phi_n(x)$ is a time-independent spatial function and $\omega_n$ is a real constant.\par
Explain why this state is called a \emph{stationary state} by explicitly discussing the probability density $|\Psi_n(x,t)|^2$ and whether it depends on time.}

\ScenarioBlock{A particle is in a 1D bound system whose (normalized) energy eigenstates are $\{\psi_n(x)\}$ with energies $\{E_n\}$. Its time-dependent state can be written as
\[
\Psi(x,t)=\sum_{n=1}^{\infty} c_n\,\psi_n(x)\,e^{-iE_n t/\hbar},
\]
with complex coefficients $c_n$ and orthonormal eigenfunctions.}

\OEBlock{Q4}{\textbf{(a)} In the expansion above, what is the physical meaning of $|c_n|^2$?}

\OEBlock{Q5}{\textbf{(b)} Using the coefficients from the expansion above, write an expression for the expectation value of the energy $\langle H\rangle$ in terms of $|c_n|^2$ and the energy eigenvalues $E_n$.}

\ScenarioBlock{Consider a particle of mass $m$ in a 1D infinite square well of width $L$:
\[
V(x)=\begin{cases}
0, & 0\le x\le L\\
\infty, & \text{otherwise.}
\end{cases}
\]
Inside the well the energy is $E>0$ and you may define $k^2=\frac{2mE}{\hbar^2}$.}

\OEBlock{Q6}{\textbf{(a)} Write down the time-independent Schr\"odinger equation in the region $0<x<L$ and give the most general form of the solution $\psi(x)$ in that region. (You may use $k^2=\frac{2mE}{\hbar^2}$.)}

\OEBlock{Q7}{\textbf{(b)} Apply the boundary conditions at $x=0$ and $x=L$ to determine the allowed values of $k$ and hence the allowed energy levels $E_n$.}

\OEBlock{Q8}{\textbf{(c)} Find the normalized stationary-state wavefunctions $\psi_n(x)$ for this well (for $0<x<L$).}

\OEBlock{Q9}{\textbf{(d)} Show that two different stationary states $\psi_n(x)$ and $\psi_m(x)$ in this well are orthogonal when $n\ne m$ by evaluating (or otherwise justifying)
\[
\int_0^L \psi_n(x)\,\psi_m(x)\,dx.
\]}

\ScenarioBlock{A particle is confined in a 1D infinite square well of width $L$. The stationary states $\{\psi_n(x)\}$ are energy eigenstates with energies $\{E_n\}$, and a general state can be expressed as a superposition of these eigenstates.}

\OEBlock{Q10}{A particle in the 1D infinite square well (width $L$) is prepared at $t=0$ in an equal superposition of the first and third stationary states:
\[
\Psi(x,0)=\frac{1}{\sqrt{2}}\big(\psi_1(x)+\psi_3(x)\big).
\]
Assuming $\psi_1$ and $\psi_3$ are normalized energy eigenstates with energies $E_1$ and $E_3$, compute the expectation value of the energy $\langle H\rangle$ at $t=0$.}

\end{document}
